\documentclass[11pt]{article}
\usepackage[margin=2.3cm]{geometry}
\usepackage{booktabs}
\usepackage{parskip}
\setlength{\parskip}{6pt}
\newcommand{\code}[1]{\texttt{\small #1}}
\title{\vspace{-1.2cm}\textbf{Does the order of questions matter?}\\[3pt]
\large Two kinds of student, two ways of choosing what to ask next}
\author{}
\date{\small 28 September 2026}
\begin{document}
\maketitle
\vspace{-1.0cm}
\noindent\rule{\textwidth}{0.4pt}

\section*{1. What this part is about}

A tutor has to decide what to ask next. There are many clever ways to
decide, but before any of them is worth building we should check something
simpler: \emph{can the order possibly matter at all?}

It turns out that depends entirely on what kind of student you imagine.
This document sets up two kinds of student and two ways of choosing a
question, crosses them, and shows that the order is nearly irrelevant for
one student and clearly important for the other.

That is a useful thing to know early. If we had only tested against the
first kind of student, every scheduler we ever build would have looked the
same as every other, and we would have concluded, wrongly, that
scheduling does not matter.

\section*{2. The curriculum is a graph}

The concepts are not a flat list. Some have to come before others. Our
practice curriculum has eight concepts in four layers:

\begin{center}
\small
\begin{tabular}{@{}ll@{}}
\toprule
\textbf{layer} & \textbf{concepts}\\
\midrule
top & \code{A1}, \code{A2} \quad (nothing comes before these)\\
second & \code{B1} needs \code{A1}; \quad \code{B2} needs \code{A1} and
\code{A2}; \quad \code{B3} needs \code{A2}\\
third & \code{C1} needs \code{B1} and \code{B2}; \quad \code{C2} needs
\code{B2} and \code{B3}\\
last & \code{D1} needs \code{C1} and \code{C2}\\
\bottomrule
\end{tabular}
\end{center}

\noindent It is shaped deliberately. \code{B2} sits under both \code{C1}
and \code{C2}, so one weak concept blocks two later ones, and you can see
it happen. And only two concepts are open at the start, so the set of
things worth asking stays small enough to read.

\section*{3. Two kinds of student}

Both students are the same underneath. Each concept is either known or
not, and each time you ask about something they do not know, there is a
chance they pick it up. The two differ in one respect only.

\textbf{Student S1 ignores the graph.} Asking about the last concept on
day one works exactly as well as asking about it at the end. The
prerequisites are drawn on the page but they do not affect this student.

\textbf{Student S2 respects the graph.} Until the prerequisites are
genuinely known, this student learns the concept far more slowly: at a
quarter of the usual rate. Asking too early is not useless, but it is
close to wasted.

S2 is the smallest change that makes the order of questions mean
something. Nothing else about the student changes.

\section*{4. Two ways of choosing a question}

\textbf{Q1} picks at random from everything the tutor does not yet believe
is mastered. It will happily ask about the last concept first.

\textbf{Q2} picks at random from the concepts whose prerequisites the
tutor \emph{believes} are already mastered. It refuses to jump ahead.

The word ``believes'' is doing real work there. The tutor cannot see what
the student actually knows, only its own running guess. So Q2 is not
following the true state of the student; it is following an estimate that
lags behind.

\section*{5. What we expected, written down first}

Before running anything, the prediction was:

\begin{center}
\small
\begin{tabular}{@{}lcc@{}}
\toprule
 & \textbf{Q1} & \textbf{Q2}\\
\midrule
S1 & fine & fine\\
S2 & bad & fine\\
\bottomrule
\end{tabular}
\end{center}

\noindent The reasoning: against S1 a question is worth the same whenever
you ask it, so nothing the scheduler does can help much. Against S2 a
question asked too early is mostly wasted, so a scheduler that avoids
asking early should win.

Writing the prediction down first matters. If the scheduler had turned out
to matter against S1 too, or not to matter against S2, the setup would not
be able to tell schedulers apart and we would have a bug rather than a
finding.

\section*{6. What actually happened}

Four hundred simulated students, forty questions each. The same student,
by seed, appears in all four boxes, so these are fair comparisons rather
than four separate groups. Scored on concepts the student \emph{truly}
knows at the end, which is the student's own state and owes nothing to
what the tutor believes.

% GENERATED by make_tables.py from sanity_outputs/. Do not edit:
% edit the experiment, rerun make_sanity.py, then make_tables.py.
\begin{center}
\small
\begin{tabular}{@{}lcccc@{}}
\toprule
 & \textbf{Q1} & \textbf{Q2} & \textbf{Q2 $-$ Q1} & \textbf{95\% CI}\\
\midrule
S1 & 6.15 & 6.00 & $-0.15$ & $\pm 0.14$\\
S2 & 4.39 & 5.86 & $+1.47$ & $\pm 0.18$\\
\midrule
\multicolumn{3}{@{}l}{interaction (S2 $-$ S1)} & $+1.62$ & $\pm 0.16$\\
\bottomrule
\end{tabular}
\end{center}


\noindent The prediction holds. Q2 is worth about one and a half concepts
to S2 and is not worth anything useful to S1.

The bottom line is the real result. It says the \emph{effect of the
scheduler depends on which student you assume}. That is the claim, not
either row on its own, and it is why testing a scheduler against only one
kind of student can be so misleading.

One technical note. Because the same student appears in both columns, the
two differences move together, and we use that: % GENERATED by make_tables.py from sanity_outputs/. Do not edit:
% edit the experiment, rerun make_sanity.py, then make_tables.py.
$r=+0.49$, and the interval would be $\pm 0.23$ rather than $\pm 0.16$ if the two were treated as independent.


\section*{7. The same result, in questions instead of concepts}

Counting concepts at a fixed budget is a bit abstract. A more natural
question is: how many questions before this student knows everything?

% GENERATED by make_tables.py from sanity_outputs/. Do not edit:
% edit the experiment, rerun make_sanity.py, then make_tables.py.
\begin{center}
\small
\begin{tabular}{@{}lccccl@{}}
\toprule
 & \textbf{Q1} & \textbf{Q2} & \textbf{mean diff} & \textbf{median diff} & \textbf{favours}\\
\midrule
S1 & 54 & 54 & $+0.3 \pm 2.3$ & $+0.0$ & unresolved at this n\\
S2 & 78 & 55 & $-22.4 \pm 4.3$ & $-17.5$ & Q2\\
\bottomrule
\end{tabular}
\end{center}


\noindent S2 gets there about twenty-three questions sooner under Q2,
which is roughly a third less work. For S1 the two are level.

Every student in every box reaches all eight concepts inside the three
hundred question budget, so nothing is cut off at the end and these are
honest medians.

Two cautions. Every box asks the same fixed number of questions, so what
differs is how quickly the student actually learns, not how much teaching
was spent. And the mean and the median are not the same number here, so
both are given rather than one standing in for the other.

\section*{8. Is it the rule, or is it the lag?}

Q1 and Q2 differ in two ways at once, and it is worth separating them.
Q2 \emph{restricts} what may be asked, and it decides readiness from a
belief that runs behind the truth. Either could be responsible.

So we ran a matched pair where both sides gate on the student's true state
and only the restriction differs. No real tutor could do this; it is a
measurement device.

% GENERATED by make_tables.py from sanity_outputs/. Do not edit:
% edit the experiment, rerun make_sanity.py, then make_tables.py.
\begin{center}
\small
\begin{tabular}{@{}lccc@{}}
\toprule
 & \textbf{unrestricted} & \textbf{restricted} & \textbf{effect}\\
\midrule
S1 & 7.480 & 7.480 & $+0.000 \pm 0.000$\\
S2 & 5.513 & 7.480 & $+1.968 \pm 0.167$\\
\bottomrule
\end{tabular}
\end{center}


\noindent For S1 the restriction is \emph{exactly} neutral. Not close to
zero, identical, for all four hundred students. That settles it: the small
negative figure for S1 in the earlier table is not the cost of restricting
what gets asked. It is the cost of deciding readiness from a belief that
lags. For S2 the restriction itself is worth about two concepts.

\section*{9. Does it hold if we run it again?}

A single run gives an interval, which tells you how well that one sample
pins the number down. It does not tell you whether the same answer would
come back from a fresh sample. For a small effect those are different
questions, so we reran on six separate groups.

% GENERATED by make_tables.py from sanity_outputs/. Do not edit:
% edit the experiment, rerun make_sanity.py, then make_tables.py.
S1 averages $+0.37 \pm 0.57$ questions of cost to Q2 over 6 independent blocks, ranging from $-0.93$ to $+1.08$, but the blocks do not agree in sign, so the direction is still unsettled; S2 averages $-23.08 \pm 1.36$ questions sooner under Q2 over 6 independent blocks, ranging from $-25.61$ to $-20.66$, and every block agrees in sign, so the direction is settled.


\noindent So the S2 result is settled and the S1 result is not. Six groups
put S1 slightly on the costly side of zero, but they disagree among
themselves.

This is worth being blunt about, because an earlier version of this work
got it wrong. We read S1 off a single group, saw a negative number, and
wrote that Q2 costs S1 nothing. The interval covered zero, so that was
asserting a direction the data never supported, and it happens to point
the opposite way from the estimate we now have. The honest statement is
that the S1 effect is too small for us to give it a sign.

\section*{10. What we are not claiming}

The effect sizes depend on the budget. At forty questions the scheduler
gap is large; run long enough for everything to be learned anyway and
every box saturates and no scheduler can differ from any other. The short
budget is a deliberate choice, not an accident.

All of this is on one small eight concept curriculum, and nothing has been
rerun on the larger one.

And S2 is a made up student. It is the simplest thing that makes
prerequisites bite, chosen for that reason, not because we have evidence
that real learners behave this way.

\section*{11. Running it again}

\code{python3 make\_sanity.py} reruns every study and writes the result
files. \code{python3 make\_tables.py} turns those into every table in this
document, so no number here is typed by hand.
\code{python3 make\_tables.py -{}-check} fails if the document and the
results ever drift apart. \code{streamlit run app.py} steps through a single run one question
at a time, with the graph on screen.

\end{document}
